\documentclass[11pt,a4paper]{article}

\usepackage[utf8]{inputenc}
\usepackage[T1]{fontenc}
\usepackage[english]{babel}
\usepackage{lmodern}
\usepackage{geometry}
\usepackage{setspace}
\usepackage{enumitem}
\usepackage{hyperref}
\usepackage{listings}
\usepackage{xcolor}

\geometry{margin=2.5cm}
\setstretch{1.12}
\setlist[itemize]{topsep=4pt,itemsep=3pt,leftmargin=18pt}
\hypersetup{colorlinks=true, linkcolor=black, urlcolor=blue}

\lstset{
  basicstyle=\ttfamily\small,
  breaklines=true,
  frame=single,
  columns=fullflexible
}

\title{Operational Guide\\\large Air Quality Assessment in Torino}
\author{[Student 1] \and [Student 2] \and [Student 3]}
\date{Academic Year 2025-2026}

\begin{document}

\maketitle

\section{Purpose}

This document explains how to reproduce the project deliverables required by the Data Management course: data acquisition, storage, transformation, quality checks, analytical queries, and dashboard execution.

\section{Repository Structure}

Main folders:

\begin{itemize}
  \item \texttt{dlt\_pipelines/}: Python ingestion pipelines
  \item \texttt{dbt\_project/}: dbt models and tests
  \item \texttt{queries/}: analytical SQL queries
  \item \texttt{dashboard\_streamlit\_v2/}: Streamlit dashboard
\end{itemize}

\section{Prerequisites}

Required software:

\begin{itemize}
  \item Python 3.13 or compatible version available in the local virtual environment
  \item Google Cloud project with BigQuery enabled
  \item dbt Core with the BigQuery adapter
  \item a valid service account key for BigQuery access
\end{itemize}

\section{Credentials and Environment Variables}

The project requires Google Cloud credentials.

Expected variables:

\begin{itemize}
  \item \texttt{GOOGLE\_APPLICATION\_CREDENTIALS}
  \item \texttt{GCP\_PROJECT\_ID}
  \item optionally \texttt{BQ\_MARTS\_DATASET} if a dataset other than \texttt{staging\_marts} is used
\end{itemize}

Recommended PowerShell setup:

\begin{lstlisting}[language={}]
$env:GOOGLE_APPLICATION_CREDENTIALS = "C:\Users\Pc\gcp-keys\dlt-loader-key.json"
$env:GCP_PROJECT_ID = "dm-torino-airquality"
\end{lstlisting}

The dashboard also supports a fallback credential path at:

\begin{lstlisting}[language={}]
%USERPROFILE%\gcp-keys\dlt-loader-key.json
\end{lstlisting}

\section{Python Environment}

From the project root, use the existing virtual environment:

\begin{lstlisting}[language={}]
\.venv\Scripts\python.exe --version
\end{lstlisting}

If needed, install project dependencies for the dashboard layer:

\begin{lstlisting}[language={}]
\.venv\Scripts\pip.exe install -r dashboard_streamlit_v2\dashboard_streamlit_v2\requirements.txt
\end{lstlisting}

\section{Data Acquisition}

Run the three dlt pipelines from the project root.

\subsection{Open-Meteo}

\begin{lstlisting}[language={}]
\.venv\Scripts\python.exe dlt_pipelines\open_meteo.py
\end{lstlisting}

\subsection{EEA DiscoData}

\begin{lstlisting}[language={}]
\.venv\Scripts\python.exe dlt_pipelines\eea_discodata.py
\end{lstlisting}

\subsection{Reverse Geocoding}

\begin{lstlisting}[language={}]
\.venv\Scripts\python.exe dlt_pipelines\station_reverse_geocoding.py
\end{lstlisting}

Expected raw outputs:

\begin{itemize}
  \item \texttt{raw.air\_quality\_statistics\_torino}
  \item \texttt{raw.weather\_torino\_hourly}
  \item \texttt{raw.station\_reverse\_geocoding}
\end{itemize}

\section{Transformation Layer with dbt}

Move into the dbt project and run transformations and tests.

\begin{lstlisting}[language={}]
Set-Location dbt_project
..\.venv\Scripts\dbt.exe run
..\.venv\Scripts\dbt.exe test
\end{lstlisting}

Expected result:

\begin{itemize}
  \item staging and mart models materialized successfully
  \item 30 dbt tests passed
\end{itemize}

Main analytical outputs:

\begin{itemize}
  \item \texttt{staging\_marts.fact\_air\_quality\_annual}
  \item \texttt{staging\_marts.mart\_station\_type\_year\_summary}
  \item \texttt{staging\_marts.mart\_air\_quality\_weather\_yearly}
  \item \texttt{staging\_marts.dim\_station\_enriched}
  \item \texttt{staging\_marts.dq\_reverse\_geocoding\_quality}
\end{itemize}

\section{Analytical Queries}

The project includes several SQL files in the \texttt{queries/} folder. For example, from PowerShell:

\begin{lstlisting}[language={}]
Get-Content -Raw .\queries\annual_means_by_station.sql | bq query --use_legacy_sql=false
\end{lstlisting}

Other available queries include:

\begin{itemize}
  \item \texttt{annual\_means\_by\_station\_type.sql}
  \item \texttt{air\_quality\_station\_year\_summary.sql}
  \item \texttt{pm10\_exceedances\_by\_station.sql}
  \item \texttt{pm10\_exceedances\_by\_station\_type.sql}
\end{itemize}

\section{Dashboard Execution}

The outer dashboard folder contains a launcher so the dashboard can be started with a simple command.

From the project root:

\begin{lstlisting}[language={}]
Set-Location dashboard_streamlit_v2
..\.venv\Scripts\streamlit.exe run app.py
\end{lstlisting}

Alternative full-path command:

\begin{lstlisting}[language={}]
\.venv\Scripts\streamlit.exe run dashboard_streamlit_v2\dashboard_streamlit_v2\app.py
\end{lstlisting}

Expected local URL:

\begin{lstlisting}[language={}]
http://localhost:8501
\end{lstlisting}

\section{Troubleshooting}

\subsection{Missing Google Credentials}

Symptom: \texttt{DefaultCredentialsError} when opening the dashboard or running pipelines.

Resolution:

\begin{itemize}
  \item verify \texttt{GOOGLE\_APPLICATION\_CREDENTIALS}
  \item verify that the service account key file exists
  \item verify \texttt{GCP\_PROJECT\_ID}
\end{itemize}

\subsection{Wrong Dashboard Path}

Symptom: \texttt{streamlit run app.py} fails from the wrong folder.

Resolution:

\begin{itemize}
  \item move to the outer \texttt{dashboard\_streamlit\_v2} folder and run the launcher
\end{itemize}

\subsection{dbt or BigQuery Errors}

Recommended checks:

\begin{itemize}
  \item confirm BigQuery billing and permissions
  \item confirm dbt profile connectivity
  \item inspect \texttt{dbt\_project/target/} and \texttt{dbt\_project/logs/}
\end{itemize}

\section{Deliverables Checklist}

Before submission, ensure that the shared folder contains:

\begin{itemize}
  \item project report in PDF or LaTeX source form
  \item project presentation in PDF or LaTeX source form
  \item operational guide
  \item all project code
  \item data or instructions to reproduce the data pipeline
\end{itemize}

This checklist is aligned with the Data Management FAQ on project documentation and submission.

\end{document}